\section{Resultados e Discussão}\label{sec:resultados}

Esta seção apresenta, para cada dimensão avaliada, o resultado obtido e a sua interpretação: a
concordância das recomendações, os defeitos que a validação revelou, a conformidade da aptidão
edáfica e sua sensibilidade às decisões do autor, a avaliação de usabilidade e as limitações.

\subsection{Concordância das recomendações}\label{sec:res_concordancia}

Os 25 casos de validação foram concordantes, com taxa de 100\% e intervalo de confiança de 95\%
de [86,7\%; 100,0\%] (\autoref{tab:concordancia}). Como o limite inferior do intervalo supera os
80\% fixados na hipótese, a hipótese foi aceita. O dado mais importante, porém, não é a taxa, e
sim o intervalo: com os 15 casos que o conjunto tinha até 10/10/2026, a taxa já era de 100\%,
mas o intervalo era de [79,6\%; 100,0\%], com o limite inferior abaixo da meta. O problema não
era de acerto, e sim de tamanho de amostra, e foi resolvido com dez casos novos.

\begin{table}[htb]
\centering
\caption{Concordância das recomendações de calagem e adubação por grupo de culturas.}
\label{tab:concordancia}
\footnotesize
\begin{tabular}{lccc}
\toprule
\textbf{Grupo} & \textbf{Casos} & \textbf{Concord.} & \textbf{IC 95\% (Wilson)} \\
\midrule
Grãos          & 12 & 12 & [75,8\%; 100\%] \\
Frutíferas     &  4 &  4 & [51,0\%; 100\%] \\
Hortaliças     &  3 &  3 & [43,9\%; 100\%] \\
Cana e tabaco  &  3 &  3 & [43,9\%; 100\%] \\
Erva-mate      &  2 &  2 & [34,2\%; 100\%] \\
Tubérculos     &  1 &  1 & [20,7\%; 100\%] \\
\midrule
\textbf{Total} & \textbf{25} & \textbf{25} & \textbf{[86,7\%; 100\%]} \\
\bottomrule
\end{tabular}
\end{table}

Dois pontos limitam a leitura desse resultado. Primeiro, o conjunto ficou abaixo dos 40 casos
previstos na proposta: cada caso exige o cálculo à mão de todas as saídas, com memória
registrada, e os casos foram escolhidos por lacuna de cobertura, e não por quantidade. Os
intervalos por grupo, largos, mostram o que isso custa: a evidência é forte para o conjunto e
para os grãos, e fraca para os grupos com um a quatro casos. Segundo, sistema e oráculo partem do
mesmo Manual. A concordância atesta que a transcrição e a implementação reproduzem os critérios
publicados, e não que esses critérios sejam agronomicamente adequados a uma área específica.

O Anexo~1 do Manual oferece uma verificação que não depende do autor: ele publica a análise de
quatro glebas e as classes de P e K resultantes para alho, milho, trigo e eucalipto. O motor do
SIRAS reproduziu as 32 classificações publicadas, sem divergência, e a tabela SMP transcrita
reproduziu as três doses de calcário que o texto do anexo calcula.

\subsection{Defeitos revelados pela validação}\label{sec:res_defeitos}

Os dez casos acrescentados para estreitar o intervalo cobriram grupos e situações que o conjunto
inicial não exercitava --- dos 15 casos originais, nove eram de calagem e um só passava pelo
algoritmo de adubação dos grãos. Dois deles divergiram do sistema, e nos dois a falha era do
sistema, não do caso. Na adubação nitrogenada, o adicional de N por rendimento esperado estava
transcrito na base para todas as culturas a que se aplica, mas nenhuma linha de código o lia: a
dose de N não variava com a expectativa de rendimento, enquanto a de P e K variava. Na classe
\emph{muito alto} de P e K, no segundo cultivo, o sistema aplicava a dose de manutenção onde o
Manual indica a de reposição, porque o texto transcrito juntava as duas numa mesma frase.

Os dois defeitos estavam no sistema desde a implementação e haviam passado por mais de cinco mil
testes automatizados sem que nenhum falhasse, porque nenhum teste combinava as entradas que os
revelam. É o resultado mais instrutivo da validação: uma taxa de 100\% sobre um conjunto
concentrado no caminho mais exercitado não media o que parecia medir. O valor do oráculo
calculado à mão está em que ele é escrito a partir do Manual, e não do código, e por isso pode
discordar do código.

\subsection{Conformidade da aptidão edáfica}\label{sec:res_aptidao}

A implementação reproduziu o CCAE nos 88 casos com classe definida, tanto na classe quanto no
fator determinante, e os quatro casos de dado ausente resultaram, como especificado, no estado
indeterminado. O resultado atesta a correção e a rastreabilidade da implementação, e não a
adequação agronômica dos critérios, que não foram validados por especialistas independentes nem em
campo.

Para medir quanto a classe depende das decisões do autor, o conjunto foi executado em quatro
configurações (\autoref{tab:sensibilidade}). A regra de rebaixamento por acúmulo, a decisão de
maior fragilidade do caderno, é também a de maior efeito: ativada, mudaria 12 dos 88 casos, quase
todos de \emph{restrita} para \emph{inapta sem correção prévia}, e dobraria essa última classe.
Todas as mudanças ocorreriam no cenário atual, nunca no potencial, porque após a correção restam
apenas os fatores permanentes. Esse efeito justifica mantê-la desativada no resultado oficial.
Os fatores permanentes pesam pouco: retirar a textura, ou a textura e a CTC, muda dois e cinco
casos, todos de \emph{apta com restrições} para \emph{apta}.

\begin{table}[htb]
\centering
\caption{Distribuição das classes de aptidão nos 88 casos, por configuração.}
\label{tab:sensibilidade}
\footnotesize
\begin{tabular}{lccccc}
\toprule
\textbf{Configuração} & \textbf{A} & \textbf{AR} & \textbf{R} & \textbf{I} & \textbf{Mudam} \\
\midrule
Sistema (fator mais limitante) & 29 & 30 & 18 & 11 & --- \\
Com rebaixamento por acúmulo   & 29 & 29 &  8 & 22 & 12 \\
Sem textura (F7)               & 31 & 28 & 18 & 11 &  2 \\
Sem CTC e textura (F5 e F7)    & 34 & 25 & 18 & 11 &  5 \\
\bottomrule
\end{tabular}

\vspace{2pt}
\raggedright\scriptsize A: apta; AR: apta com restrições; R: restrita; I: inapta sem correção prévia.
\end{table}

A comparação entre cenários é a parte da aptidão que integra os dois módulos do sistema. No caso
da Gleba~4 com milho (\autoref{fig:aptidao}), a área é \emph{restrita} hoje, limitada pelo
potássio, e passa a \emph{apta} após a correção recomendada pelo próprio SIRAS, quando a única
limitação remanescente é a CTC, que não se corrige com insumos. O laudo deixa assim de dizer
apenas quanto aplicar e passa a dizer o que a aplicação muda.

\subsection{Avaliação de usabilidade}\label{sec:res_usabilidade}
Participaram oito pessoas: cinco técnicos agrícolas, dois estudantes de agronomia e um
engenheiro agrônomo. Quatro têm mais de dez anos de trabalho com análise de solo ou recomendação
de adubação; os demais se distribuem entre seis e dez anos (um), dois a cinco anos (um), menos de
dois anos (um) e ainda sem atuação profissional (um). Hoje, cinco fazem a recomendação à mão, com
o Manual, quatro usam o sistema do laboratório, dois um programa de cooperativa ou revenda e um
uma planilha própria --- a pergunta admitia mais de uma resposta. Todos se declararam
confortáveis (quatro) ou muito confortáveis (quatro) com ferramentas digitais, nenhum conhecia o
SIRAS antes do teste e um deles (P7) já tinha usado o FertFácil.

Os oito concluíram sozinhos o cenário fechado do segundo roteiro, sem intervenção do pesquisador.
No questionário, das 80 avaliações possíveis (oito participantes, dez heurísticas), 79 receberam
nota; em uma o participante indicou não ter passado pela situação. Foram 73 notas 0 (\emph{não é
problema}) e seis notas 1 (\emph{problema estético}); nenhuma heurística recebeu nota 2 ou maior
(\autoref{tab:heuristicas}). As seis notas 1 vieram de três participantes, e quatro delas de um só
(P5). A mediana foi 0 em todas as heurísticas, e a ordenação por gravidade que o protocolo previa
não se aplica: as diferenças entre heurísticas são de uma nota, de um participante.

\begin{table}[htb]
\centering
\caption{Gravidade atribuída pelos participantes a cada heurística ($n=8$).}
\label{tab:heuristicas}
\footnotesize
\begin{tabular}{llccc}
\toprule
\# & Heurística & Nota 0 & Nota 1 & Não observou \\
\midrule
H1  & Visibilidade do status do sistema      & 7 & 1 & -- \\
H2  & Correspondência com o mundo real       & 7 & 1 & -- \\
H3  & Controle e liberdade do usuário        & 8 & -- & -- \\
H4  & Consistência e padrões                 & 8 & -- & -- \\
H5  & Prevenção de erros                     & 7 & 1 & -- \\
H6  & Reconhecimento em vez de memorização   & 8 & -- & -- \\
H7  & Flexibilidade e eficiência de uso      & 8 & -- & -- \\
H8  & Estética e design minimalista          & 7 & 1 & -- \\
H9  & Ajuda para reconhecer e corrigir erros & 6 & 1 & 1 \\
H10 & Ajuda e documentação                   & 7 & 1 & -- \\
\bottomrule
\end{tabular}

\vspace{2pt}
\raggedright\scriptsize Escala de 0 (não é problema) a 4 (catástrofe). Nenhuma heurística
recebeu nota 2, 3 ou 4.
\end{table}

Os relatos em texto são mais informativos que as notas. Dois descrevem a mesma falha, de
prevenção e recuperação de erro no formulário: P5 observou que a tecla Enter gera o laudo mesmo
com dados incompletos e sugeriu uma confirmação; P7, que de início não informou um campo e recebeu
apenas uma mensagem de erro, sem indicação do que faltava, e pediu que os campos pendentes fossem
destacados. Os dois pontos foram corrigidos depois da coleta, fora da versão avaliada: a tecla
Enter passou a levar ao campo seguinte, e a recusa de um valor passou a nomear e destacar o campo.
P2 relatou ter conseguido informar um só tipo de adubo da revenda, embora o formulário aceite até
três fórmulas. Dois participantes apontaram o contexto de uso: P1 lembrou que o teste não reproduz
o dia a dia com internet lenta, e P5 pediu uma versão que funcione sem conexão e fora do
computador do escritório.

Nas perguntas abertas, os oito disseram que usariam a ferramenta no trabalho, com justificativas
centradas em tempo e precisão --- \emph{``ter o laudo completo em menos de 10 min''} (P1),
\emph{``evita erros, e possui linguagem agronômica de fácil entendimento''} (P4). O participante
que já conhecia o FertFácil destacou justamente o que o diferencia: \emph{``as informações de onde
saiu os números achei muito bom''} (P7). Sobre a recomendação, os oito a consideraram confiável e
a recomendariam ao produtor; P1 ressalvou que conferiria os valores à mão, e P4 citou a adaptação
da recomendação às fórmulas disponíveis nas revendas da região. Seis responderam que não falta
nada para a sua rotina; as duas exceções são as já citadas: a versão sem conexão (P5) e o
destaque dos campos pendentes (P7).

As notas devem ser lidas como um teto, e não como prova de ausência de problemas. A amostra é
pequena e por conveniência, o questionário foi respondido logo após uma sessão acompanhada pelo
pesquisador, e esses fatores favorecem avaliações positivas. O que o teste sustenta é mais
modesto e mais útil: profissionais do público-alvo, inclusive um sem experiência e um com mais de
dez anos, concluíram sozinhos um cenário de recomendação completo e não encontraram obstáculo que
julgassem mais que estético; e os relatos localizaram um ponto concreto a melhorar, a
sinalização de dados faltantes no formulário.

\subsection{Comparação com ferramenta existente}\label{sec:res_comparacao}
O SIRAS foi posicionado frente ao FertFácil \citep{Tiecher2016, FertFacil2026}, a ferramenta de
recomendação de calagem e adubação mais difundida no RS, adotada pela Emater/RS. A
\autoref{tab:fertfacil} resume a comparação funcional, feita a partir da documentação pública e de
um laudo de exemplo do FertFácil.

\begin{table}[htb]
\centering
\caption{Comparação funcional entre o FertFácil e o SIRAS.}
\label{tab:fertfacil}
\footnotesize
\begin{tabularx}{\columnwidth}{lXX}
\toprule
 & \textbf{FertFácil} & \textbf{SIRAS} \\
\midrule
Base técnica & Manuais RS/SC (2016) e PR (2019) & Manual RS/SC (2016) \\
Escopo & todas as culturas dos dois manuais & 61 culturas e grupos \\
Calagem e adubação & sim & sim \\
Indicação do adubo & uma fórmula, com dose e complemento & fórmula ideal, catálogo e mistura de simples \\
Aptidão edáfica & não & atual e potencial \\
Origem de cada valor & não exibida & regra e tabela do Manual \\
Laudo & sintético, imprimível & essencial e anexo explicativo, imprimível \\
Execução & on-line & local, sem internet \\
Maturidade & em uso desde 2015 & protótipo acadêmico \\
\bottomrule
\end{tabularx}
\end{table}

As duas ferramentas convergem no núcleo --- recomendar calagem e adubação pelo Manual RS/SC --- e
divergem no que entregam em torno dele. O FertFácil cobre mais: dois manuais completos, e não um
recorte de 61 culturas, além de um módulo que compara 11 manuais para a soja. É também uma
ferramenta madura, em uso profissional há mais de uma década. O SIRAS acrescenta três coisas que
o FertFácil não oferece: a estimativa de aptidão edáfica, com o efeito da correção sobre ela; a
indicação da origem de cada número do laudo; e um laudo que, além da recomendação, explica a
interpretação de cada atributo e o caminho até cada dose. O FertFácil entrega a receita; o SIRAS
entrega a receita com a justificativa. As diferenças de escopo e de maturidade são esperadas entre
um produto consolidado e um trabalho de conclusão de curso, e a integração entre recomendação e
aptidão, ausente nas ferramentas analisadas, é a contribuição que este trabalho se propôs a fazer.

A proposta previa também submeter um subconjunto dos casos de validação ao FertFácil e comparar as
saídas, e a mesma restrição impediu uma comparação de usabilidade entre as duas ferramentas. Nenhuma
das duas foi realizada, porque o acesso ao sistema, que exige cadastro, não pôde ser obtido no
período da avaliação. A perda é limitada: as duas ferramentas implementam o mesmo Manual,
e a correção numérica do SIRAS já foi medida contra o oráculo calculado à mão. A comparação
numérica permanece como trabalho futuro, útil sobretudo para localizar pontos de leitura
divergente do Manual entre as duas implementações.

\subsection{Limitações}\label{sec:res_limitacoes}

A aptidão é exclusivamente edáfica: não considera relevo, clima, drenagem nem erosão, e seus
critérios, embora ancorados no Manual e no sistema oficial, foram definidos pelo autor e não
foram validados por especialistas independentes nem em campo. O oráculo das recomendações deriva
do mesmo Manual que o sistema implementa, o que limita a validação à fidelidade aos critérios
publicados. O conjunto de casos é menor que o previsto e desigual entre os grupos. A avaliação de
usabilidade usa amostra pequena e por conveniência, com sessões acompanhadas pelo pesquisador; as
notas de severidade exprimem percepção dos usuários, não gravidade técnica, e a pergunta sobre a
recomendação reuniu o juízo sobre a ferramenta e o juízo sobre a recomendação agronômica. Por fim, o SIRAS é ferramenta de apoio à decisão e não
substitui a responsabilidade técnica de profissional habilitado.
